\chapter{Fundamentação Teórica e Trabalhos Relacionados}
\label{cap:cap2}

Para fundamentar a proposta deste trabalho e identificar as lacunas no estado da arte, foi realizada uma mini-revisão sistemática da literatura. O protocolo de busca utilizou as bases de dados Google Scholar e IEEE Xplore, considerando publicações entre os anos de 2019 e 2024. As strings de busca utilizadas combinaram os termos: "IoT dam monitoring", "Computer Vision in structural health monitoring" e "Low-power wide-area networks for water management".

De forma complementar, foi utilizado o modelo de linguagem Gemini Pro como ferramenta de apoio à exploração bibliográfica. A estratégia consistiu na elaboração de prompts descritivos sobre a temática do trabalho, solicitando especificamente a busca por estudos recentes apresentados em conferências de alta relevância. Para garantir a integridade dos dados, cada trabalho sugerido pelo modelo foi submetido a conferência humana, confrontando as citações e resumos gerados com o conteúdo dos textos originais.

Foram selecionados trabalhos que atendessem aos seguintes critérios de inclusão: (1) propostas de arquiteturas de IoT para monitoramento hídrico ou estrutural; (2) uso de processamento na borda (Edge Computing); (3) uso de métodos de Visão Computacional e IA para monitoramento de estruturas; e (4) aplicação em cenários com recursos limitados. A análise da literatura buscou responder à seguinte questão norteadora: "Como as soluções atuais de monitoramento de barragens lidam com as restrições de comunicação e a dependência da computação em nuvem, principalmente quando exigem processamento de imagens?".

A seguir, apresentam-se os trabalhos mais relevantes identificados e como eles se relacionam com a abordagem proposta.

A arquitetura de referência para Smart Water Cities proposta por \citeonline{oberascher2022} vincula a escolha da infraestrutura de comunicação às demandas de resolução espacial e temporal, destacando que "o conceito de IoT permite novas abordagens para monitoramento e controle da infraestrutura hídrica urbana [...], mesmo em estruturas remotas ou subterrâneas". Embora o framework original associe o monitoramento por imagens a redes de alta capacidade, este trabalho inova ao integrar a Visão Computacional como um sensor IoT na borda (Edge Computing). Ao processar o fluxo visual localmente e extrair apenas metadados de anomalias (fissuras), a solução proposta rompe com a dependência de banda larga, viabilizando o uso de imagens como forma de sensoriamento para barragens localizadas em áreas remotas, sem a necessidade de deslocamento físico constante.

\citeonline{bastos2023} apresenta um sistema de aquisição de dados universal para barragens utilizando IoT, focando na equalização de piezômetros. O presente trabalho avança sobre essa premissa ao incorporar não apenas sensores escalares, mas sensores visuais (câmeras) interpretados por algoritmos de Aprendizado Profundo (Deep Learning).

Em uma revisão crítica recente sobre o estado da arte do monitoramento de saúde estrutural (SHM - Structural Health Monitoring) em barragens, \citeonline{khan2024} enfatizam que a integração de IoT e Inteligência Artificial revolucionou a área ao permitir a coleta e análise de dados em tempo real, superando as limitações de frequência irregular e subjetividade das inspeções visuais tradicionais. Os autores identificam, contudo, que desafios significativos persistem: notadamente a necessidade de sistemas de comunicação robustos e a garantia de confiabilidade dos dados em ambientes severos. O presente trabalho estende a revisão de \citeonline{khan2024} ao abordar o monitoramento visual contínuo, ainda pouco explorado no SHM de barragens, e busca solucionar um dos principais entraves práticos para sua execução: a conectividade limitada em regiões remotas. Por meio de uma arquitetura de Edge Computing, o processamento de imagens ocorre localmente. A escolha pela arquitetura de Edge Computing justifica-se pela latência e limitações de banda em áreas rurais. Diferente de abordagens que dependem inteiramente da nuvem, o processamento na borda permite inferência em tempo real e redução de tráfego de rede, uma estratégia essencial para infraestruturas críticas isoladas.

Paralelamente, \citeonline{kamal2022} reforçam que a transição para sistemas sem fio (wireless) e sensores miniaturizados é a solução definitiva para os problemas de custo proibitivo e tempo de instalação dos sistemas cabeados tradicionais. O presente trabalho também se alinha a essa perspectiva, já que se torna uma alternativa à utilização de extensômetros (uma única câmera pode substituir a necessidade de instalação de N sensores). No experimento desenvolvido e apresentado adiante, a tecnologia sem fio LoRa foi utilizada para a transmissão local dos dados entre os dispositivos telemétricos.

No âmbito da viabilidade técnica do processamento na borda, \citeonline{kim2023} propõem uma arquitetura de inferência distribuída para detecção de fissuras em ambientes IoT. Os autores demonstram que, ao segmentar o processamento de redes neurais profundas (baseadas em EfficientNet e U-Net) entre o dispositivo de captura e um servidor de borda, é possível reduzir drasticamente o tempo de inferência sem sacrificar a precisão das detecções. Apesar de os autores terem utilizado uma estratégia de particionamento de carga, a abordagem utilizada por eles corrobora a hipótese deste trabalho, de que dispositivos com recursos limitados podem ser integrados a arquiteturas de Edge Computing para viabilizar tarefas complexas de visão computacional.

Em uma aplicação direta ao contexto de infraestrutura hídrica, \citeonline{srinivas2024} desenvolveram um sistema de automação para barragens que integra sensores IoT (nível, chuva) com detecção de falhas visuais baseada no algoritmo YOLOv5. O estudo destaca a eficiência da integração entre hardware embarcado e modelos de Deep Learning para a identificação de defeitos estruturais em tempo real. Entretanto, diferentemente da proposta de Srinivas et al., que foca na automação de comportas, o presente trabalho expande a aplicação da IA para a geolocalização precisa das anomalias em um modelo tridimensional (gêmeo digital), oferecendo uma ferramenta de gestão visual mais robusta para a tomada de decisão remota.
